\section{Related Work}
\label{sec:related}

\paragraph{Critic-free RL for language models and agents.}
PPO~\cite{schulman2017ppo} with a learned value function~\cite{schulman2015gae} was the original
workhorse of RL from human feedback~\cite{ouyang2022instructgpt}. GRPO~\cite{shao2024deepseekmath}
removed the critic by standardizing each sampled response's reward against a group of responses to
the same prompt, and with its successors (REINFORCE-style
estimators~\cite{ahmadian2024back}, DAPO's clipping and sampling changes~\cite{yu2025dapo}) it became
the default for reasoning models~\cite{guo2025deepseekr1}. Carrying it to multi-turn agents required
new infrastructure (HybridFlow/verl~\cite{sheng2025hybridflow}, RAGEN~\cite{wang2025ragen}); an
alternative line keeps a critic but makes it hierarchical across turns~\cite{zhou2024archer}.

\paragraph{Step-level credit without a critic.}
GiGPO~\cite{feng2025gigpo} is, to our knowledge, the first critic-free step-level advantage for LLM
agents. It hashes environment states across the group's rollouts into \emph{anchor states} and
standardizes the discounted return of each visit against the other visits to the same state, then
adds this step advantage to the episode one. GACA~\cite{liang2026gaca} argues that the useful
granularity is a property of the state and mixes the two levels with a per-step weight $\lambda_t$
driven by a criticality score, the sampled action's mean token NLL normalized within its trajectory;
it proves that a state-adaptive weight dominates any constant one and that the optimal weight grows
with the anchor-group size. Its Appendix B.4 proposes, but does not train, a leave-one-out score that
averages the NLL of the other visits to the same anchor. Both methods stop at the environment step: all
tokens of a step's response share one advantage. Process reward models~\cite{lightman2024verify,
wang2024mathshepherd} and Monte-Carlo value estimates from intermediate states~\cite{kazemnejad2024vineppo}
offer finer credit at the cost of a learned verifier or many extra rollouts; reward redistribution
methods such as RUDDER~\cite{arjona2019rudder} reshape returns, which reward-shaping theory
constrains~\cite{ng1999shaping}. Our method adds no model and no rollouts.

\paragraph{Token-level signals in RL for reasoning.}
Wang~\etal~\cite{wang2025beyond} show that in single-turn reasoning a small minority of high-entropy
``forking'' tokens drives most of the improvement, and that restricting the policy gradient to them
helps; Cui~\etal~\cite{cui2025entropy} study how RL collapses policy entropy. GACA also scores steps by
uncertainty. We test whether uncertainty locates the \emph{decision} inside an agent's reasoning and
find that it largely does not (\cref{sec:measure}).

\paragraph{Faithfulness and causal importance of reasoning.}
Lanham~\etal~\cite{lanham2023faithfulness} test whether chain of thought is used by truncating it and
forcing an early answer. Bogdan~\etal~\cite{bogdan2025anchors} locate ``thought anchors'', the
sentences that shape a reasoning trace, by resampling the trace from each sentence and, as a white-box
alternative, by inspecting attention between sentences. Yu~\etal~\cite{yu2026cir} formalize the
truncate-and-force measurement as Causal Importance of Reasoning,
$\mathrm{CIR} = \frac{1}{T}\sum_k \mathrm{JS}\big(\mathrm{Bern}(p_k)\,\|\,\mathrm{Bern}(p_T)\big)$
with $p_k$ the probability of the final answer after the first $k$ reasoning tokens, show that outcome
rewards lower it on many tasks, and restore it with a CIR bonus in the reward. Our commitment probe is
this measurement. What we add is its use \emph{inside the advantage}: the probe's per-segment rise
decides which tokens of a step receive that step's credit, in a multi-turn agent, instead of a scalar
bonus shared by all tokens. We also check it against resampling and against attention, the two
signals Bogdan~\etal~compare.
